\section{A complete finite arithmetic test below a height}
\label{sec:fd-arithmetic}
The even trace powers of Suzuki's compressions give a finite arithmetic certificate for every bounded packet that consists of simple zeros.
Write \(\Xi(z)=\xi(1/2+iz)\), and let \(P_N\) be the orthogonal
projection in the Weil Hilbert space onto the mean-zero dyadic
piecewise-linear space \(V_N\) supported on \([-N,N]\).
Let \(C_N=P_NCP_N\), where \(C\) has the single eigenvalue
\(1/\gamma\) at each distinct signed ordinate \(\gamma\).
Its coordinate norm has weight \(m_\gamma\).
For \(k\ge1\) define
\[
 S_k=-2k[z^{2k}]\log\frac{\Xi(z)}{\Xi(0)},\qquad
 D_{N,k}=S_k-\operatorname{tr}(C_N^{2k}),\qquad
 H_{N,k}=\left(\frac2{D_{N,k}}\right)^{1/(2k)}.
\]
For a height \(T>0\), let
\[
 E(T)=\sum_{0<\gamma\le T}(m_\gamma-1).
\]
This is the multiplicity excess in the positive packet. It is an
integer and vanishes exactly when that packet is simple.

\begin{theorem}[Finite arithmetic certificates and the first multiple zero]
\label{thm:fd-arithmetic-complete}
For every \(N,k\ge1\), the number \(D_{N,k}\) is positive and
\[
 D_{N,k}\ge\frac{2E(T)}{T^{2k}}\qquad(T>0).
\]
For each fixed \(k\), it decreases with \(N\) and
\[
 \lim_{N\to\infty}D_{N,k}
   =2\sum_{\gamma>0}\frac{m_\gamma-1}{\gamma^{2k}}.
\]
If \(\gamma_*\) denotes the least positive multiple ordinate,
with value \(+\infty\) when there is none, then
\[
 \sup_{N,k\ge1}H_{N,k}=\gamma_*.
\]
Consequently, for each fixed \(T>0\),
\[
 \begin{split}
 &\text{every zero with }0<\Im\rho\le T\text{ is simple}\\
 &\hspace{16mm}\Longleftrightarrow\
 \text{there are finite }N,k\text{ with }D_{N,k}<2/T^{2k}.
 \end{split}
\]
The exact relation to the Jensen limit is
\[
 \lim_{a\downarrow0}a\,j_*(-a)
   =\frac12\left(\sup_{N,k\ge1}H_{N,k}\right)^2.
\]
\end{theorem}

\begin{proof}
The paired Hadamard product gives
\[
 S_k=2\sum_{\gamma>0}\frac{m_\gamma}{\gamma^{2k}}.
\]
The trace of \(C^{2k}\) is
\(2\sum_{\gamma>0}\gamma^{-2k}\), since each distinct signed
coordinate is one-dimensional. Compression decreases each
singular value: the approximation-number characterization gives
\(s_j(P_NCP_N)\le s_j(C)\). Summing their \(2k\)-th powers
therefore proves
\[
 D_{N,k}\ge
 2\sum_{\gamma>0}(m_\gamma-1)\gamma^{-2k}
 \ge 2E(T)T^{-2k}.
\]
The first inequality is strict before subtracting the excess:
\(C_N\) has finite rank and \(C\) has infinitely many nonzero
singular values. Thus \(D_{N,k}>0\).

The spaces \(V_N\) are nested and dense in the Weil space, as
proved in the compression construction. The min--max principle
makes each positive eigenvalue of \(C_N\) nondecreasing with \(N\);
applied to \(-C_N\), it makes the absolute values of the negative
eigenvalues nondecreasing. Adding zero eigenvalues when necessary,
their even-power sum is nondecreasing. This proves monotonicity
of \(D_{N,k}\).

Suzuki's compression argument gives \(C_N\to C\) in
Hilbert--Schmidt norm. To pass to higher powers, use
\[
 A^{2k}-B^{2k}
 =\sum_{j=0}^{2k-1}A^j(A-B)B^{2k-1-j}.
\]
In each summand, one of the powers besides \(A-B\) contains
at least one factor. Put that factor in Hilbert--Schmidt norm,
put \(A-B\) in the same norm, and bound the remaining powers
in operator norm. The ideal inequality
\(\|UV\|_1\le\|U\|_2\|V\|_2\) proves convergence in trace norm,
since the operator and Hilbert--Schmidt norms remain bounded.
The asserted limit follows.

For the supremum, list the numbers \(1/\gamma\) belonging to
multiple positive ordinates, each repeated \(m_\gamma-1\) times,
and denote this nonnegative sequence by \(v\). It belongs to
\(\ell^2\). Monotonicity in \(N\) gives
\[
 \sup_N H_{N,k}=\|v\|_{2k}^{-1},
\]
with the right side infinite for the empty sequence.
For a nonempty sequence, its largest entry is
\(v_{\max}=1/\gamma_*\), and
\[
 v_{\max}^{2k}\le\sum_jv_j^{2k}
 \le v_{\max}^{2k-2}\sum_jv_j^2.
\]
Taking \(2k\)-th roots shows
\(\|v\|_{2k}\to v_{\max}\). Moreover
\(\|v\|_{2k}\ge v_{\max}\) for every \(k\). These two facts prove
the supremum identity, including the empty case.

A strict defect bound makes the nonnegative integer \(E(T)\)
less than one and hence zero. Conversely, when \(E(T)=0\),
discreteness of the zeros gives \(\gamma_*>T\); the supremum
identity then supplies a finite pair with \(H_{N,k}>T\).
Finally \(x_*=2\gamma_*\) in Theorem~\ref{thm:fd-limits}, so
its Jensen formula is exactly the last displayed identity.
\end{proof}

This theorem supplies an exact relation among the spectral finite
data, bounded-height simplicity, and the Jensen detecting constant.
At every fixed height the strict test has a terminating search
precisely for a simple packet: enumerate \(N,k\) and increase
the certified precision until a strict inequality is found.
The assertion that this search terminates for every height is
the all-simple question.

\subsection{Operator control and the size of a finite certificate}

The same defect controls the operator approximation, without further
matrix entries. All operator norms below are for the Weil Hilbert norm.

\begin{proposition}[An operator bound from an even trace defect]
\label{prop:fd-even-operator}
For every \(N,k\ge1\),
\[
 \|C-C_N\|\le\frac2{\sqrt3}\,D_{N,k}^{1/(2k)}.
\]
The constant \(2/\sqrt3\) is sharp in the elementary selfadjoint
block estimate used in this proof.
\end{proposition}
\begin{proof}
Write \(P=P_N\), \(Q=I-P\), and let \(v_j\) be an orthonormal
eigenbasis for the finite compression on \(P\mathcal H\).
Scalar convexity, applied to the spectral measure of \(C\) at
\(v_j\), gives
\[
 |\langle Cv_j,v_j\rangle|^{2k}
       \le\langle |C|^{2k}v_j,v_j\rangle.
\]
After summing, and using \(S_k\ge\operatorname{tr}|C|^{2k}\),
we obtain
\[
 D_{N,k}\ge\operatorname{tr}(Q|C|^{2k}Q).
\]
For a unit vector \(v\in Q\mathcal H\), the power-mean
inequality for its spectral measure gives
\[
 \|Cv\|^{2k}\le\langle |C|^{2k}v,v\rangle
       \le\operatorname{tr}(Q|C|^{2k}Q)\le D_{N,k}.
\]
Thus \(\delta=\|CQ\|\le D_{N,k}^{1/(2k)}\).

Relative to \(P\mathcal H\oplus Q\mathcal H\), write
\[
 C-C_N=\begin{pmatrix}0&B^*\\B&D\end{pmatrix}.
\]
The column bound is \(BB^*+D^2\le\delta^2I\).
For a unit vector \((x,y)\), put \(a=\|x\|\),
\(b=\|y\|\), and \(w=\|Dy\|\). Then
\begin{align*}
 |\langle(C-C_N)(x,y),(x,y)\rangle|
 &\le2a\sqrt{\delta^2b^2-w^2}+bw\\
 &\le\delta\sqrt{4a^2b^2+b^4}
   =\delta\sqrt{4b^2-3b^4}
   \le\frac{2\delta}{\sqrt3}.
\end{align*}
The second inequality is the ordinary two-dimensional
Cauchy--Schwarz inequality. The last polynomial attains its
maximum at \(b^2=2/3\). Taking the supremum proves the assertion.
Sharpness of this block estimate follows from the scalar blocks
\(B=\delta\sqrt{2/3}\), \(D=\delta/\sqrt3\), whose omitted
column has norm \(\delta\) and whose resulting block matrix has
norm \(2\delta/\sqrt3\).
\end{proof}

Let \(\gamma_1\le\gamma_2\le\cdots\) now list the positive
ordinates with multiplicity. This counted list is used only in the
following exact size comparison.

\begin{proposition}[Dimension and moment order]
\label{prop:fd-certificate-size}
Put \(d_N=\dim V_N=2r_N\). For every \(k\ge1\),
\[
 D_{N,k}\ge2\sum_{j>r_N}\gamma_j^{-2k}.
\]
Consequently a strict certificate \(D_{N,k}<2/T^{2k}\)
requires \(d_N\ge2n(T)\), where \(n(T)\) counts positive
zeros through \(T\) with multiplicity. For each fixed \(N\),
\[
 \lim_{k\to\infty}H_{N,k}=\gamma_1.
\]
\end{proposition}
\begin{proof}
The rank is at most \(d_N\), and compression decreases every
singular value. Hence \(\operatorname{tr}(C_N^{2k})\) is at
most the sum of the largest \(d_N\) even powers of the singular
values of \(C\). The counted signed ordinate list contains each
such singular value at least as many times as the distinct
coordinate list. Its largest \(2r_N\) even powers sum to
\(2\sum_{j\le r_N}\gamma_j^{-2k}\). Subtraction from \(S_k\)
proves the tail bound. If \(n(T)>r_N\), its first omitted
term is at least \(2/T^{2k}\), contradicting the strict
certificate. This proves the dimension requirement.

Write \(\lambda_1=1/\gamma_1=\|C\|\) and
\(\beta_N=\|C_N\|\). We claim \(\beta_N<\lambda_1\).
Equality would give a unit vector in \(V_N\) attaining an
extreme Rayleigh value of \(C\). Positivity of
\(\lambda_1I\mp C\) would make it a \(C\)-eigenvector with
eigenvalue \(\pm\lambda_1\). The compact-primitive identity
then gives, as equality in the Weil space,
\[
 i\int_{-N}^{x}f(y)\,dy=\pm\lambda_1 f(x).
\]
The injective compact-support embedding makes this an equality
of ordinary functions. Differentiating gives a first-order
constant-coefficient equation, with \(f(-N)=0\), and therefore
\(f=0\). This contradicts its unit norm and proves the claim.

The summable counted reciprocal-square sequence gives
\(S_k^{1/(2k)}\to\lambda_1\). Moreover,
\[
 0\le\frac{\operatorname{tr}(C_N^{2k})}{S_k}
 \le\frac{d_N}{2}
       \left(\frac{\beta_N}{\lambda_1}\right)^{2k}
 \longrightarrow0.
\]
It follows that \(D_{N,k}^{1/(2k)}\to\lambda_1\), which is
exactly the last assertion.
\end{proof}

The dimension requirement and the fixed-space limit quantify two
constraints on the actual finite search. Increasing the moment
order alone in one space eventually returns its height threshold
to the first ordinate. The complete family in
Theorem~\ref{thm:fd-arithmetic-complete} permits the space and the
moment order to vary together.

\subsection{Evaluation using the arithmetic side of Weil's formula}

We now evaluate these certificates without inserting zero ordinates.
For \(h=f*\widetilde g\), where
\(\widetilde g(x)=\overline{g(-x)}\), the explicit formula is
\begin{align}
 q(f,g)
 ={}&\int_{\mathbb R}h(x)(e^{x/2}+e^{-x/2})\,dx
       -(\log\pi+\gamma_{\!E})h(0)\notag\\
 &+\int_0^\infty
   \frac{2e^{-2x}h(0)-e^{-x/2}(h(x)+h(-x))}
        {1-e^{-2x}}\,dx\notag\\
 &-\sum_{n\ge2}\frac{\Lambda(n)}{\sqrt n}
             \bigl(h(\log n)+h(-\log n)\bigr).
 \label{eq:fd-arithmetic-weil}
\end{align}
Here \(\gamma_{\!E}\) is Euler's constant and \(\Lambda(n)\)
is the von Mangoldt function. This normalization follows from
\cite[(3.2)--(3.7), pp.~5--6]{ccm2025} by \(u=e^x\).
Their constant is \(-(\log(4\pi)+\gamma_{\!E})h(0)\), with
\(2e^{-x}h(0)\) in the numerator. The identity
\[
 2\int_0^\infty\frac{e^{-2x}-e^{-x}}{1-e^{-2x}}\,dx=-\log4
\]
proves equality with~\eqref{eq:fd-arithmetic-weil}.
Only prime powers \(n<e^{2N}\) enter on the support in question.

Put \(h_N=2^{-N}\) and
\[
 t_j(x)=\max\{1-|x-jh_N|/h_N,0\},\qquad
 b_j=t_j-t_0,
\]
where \(-N/h_N<j<N/h_N\) and \(j\ne0\).
These \(2N2^N-2\) functions form the stated basis of \(V_N\):
all interior hats have the same integral \(h_N\), so their
coefficient sum is exactly the mean-zero condition.
The correlation of the central hat is
\[
 b(x)=(t_0*t_0)(x)
 =\frac1{6h_N^2}\sum_{\ell=0}^4
       (-1)^\ell\binom4\ell
             (x+(2-\ell)h_N)_+^3.
\]
It can be checked by integration on the two intervals:
\[
 b(x)=
 \begin{cases}
 2h_N/3-x^2/h_N+|x|^3/(2h_N^2),&|x|\le h_N,\\
 (2h_N-|x|)^3/(6h_N^2),&h_N\le|x|\le2h_N,\\
 0,&|x|\ge2h_N .
 \end{cases}
\]
Let \(W\) denote the linear functional in
\eqref{eq:fd-arithmetic-weil}, and define real kernels
\[
 g(d)=W(b(\,\cdot-d)),\qquad
 a(d)=W\left(\int_0^d b(\,\cdot-u)\,du\right).
\]
The first kernel is even and the second odd. In the Fourier
variable,
\[
 a(d)=\sum_{\gamma\in\Gamma}
      m_\gamma|\widehat t_0(\gamma)|^2
                   \frac{\sin(\gamma d)}{\gamma}.
\]
The identity follows by integrating the absolutely convergent
correlation series; equivalently it follows from the compact
primitive, after taking the basis differences.
Consequently the Gram and primitive matrices are
\[
 \begin{split}
 G_{ij}&=g((j-i)h_N)-g(jh_N)-g(ih_N)+g(0),\\
 B_{ij}&=a((j-i)h_N)-a(jh_N)+a(ih_N),\\
 H_{ij}&=q\left(i\int_{-N}^{\,\cdot}b_j,b_i\right)=-iB_{ij}.
 \end{split}
\]
Thus \(G\) is real symmetric and \(B\) real antisymmetric, and
\[
 \operatorname{tr}(C_N^{2k})
       =(-1)^k\operatorname{tr}\bigl((G^{-1}B)^{2k}\bigr).
 \label{eq:fd-arithmetic-trace}
\]
This evaluates the invariant trace without a matrix square root.

\subsection{Certified examples}

The following are strict rational height certificates from
\eqref{eq:fd-arithmetic-weil}. Each displayed upper bound \(U\)
encloses \(D_{N,k}\), and each row satisfies the exact rational
inequality \(U<2/T^{2k}\).
\[
\begin{array}{c|r|r|c|c|c}
 N&\dim V_N&\text{prime powers}&k&U&T\\ \hline
 1&2&5&9&3.946530144\,10^{-21}&14.134\\
 2&14&24&3&2.539105111\,10^{-8}&20.704\\
 3&46&98&3&5.419648073\,10^{-9}&26.782\\
 4&126&465&2&7.880330238\,10^{-7}&39.913\\
 5&318&2532&2&1.318023575\,10^{-7}&62.413\\
 6&766&15040&3&2.309270917\,10^{-12}&97.632
\end{array}
\]
In particular, this finite arithmetic calculation certifies
simplicity of every nontrivial zero with
\(0<\Im\rho\le97.632\).
The largest prime power in that calculation is \(162751\).
Proposition~\ref{prop:fd-even-operator} also gives the strict
operator estimate \(\|C-C_6\|<0.0133\): raising this bound
to the sixth power reduces its verification to the exact rational
inequality \(64(2.309270917\cdot10^{-12})<27(0.0133)^6\).

Here is the numerical error control. On each dyadic interval,
\(h(x)+h(-x)\) is a rational polynomial of degree at most four.
The polar integrals are polynomial-exponential antiderivatives.
The remaining integrals have the form
\[
 I_{\ell,r}=\int_{\ell h_N}^{(\ell+1)h_N}
       \frac{x^r e^{-x/2}}{1-e^{-2x}}\,dx,\qquad 0\le r\le4.
\]
At the origin the constant term is cancelled before evaluation.
Writing
\(\operatorname{exprel}(z)=(e^z-1)/z={}_1F_1(1;2;z)\),
the regular integrand for \(r\ge1\) is
\[
 \frac{x^{r-1}e^{-x/2}}{2\operatorname{exprel}(-2x)}.
\]
The constant-term integrand is
\[
 -\frac32 e^{-x/2}
       \frac{\operatorname{exprel}(-3x/2)}
            {\operatorname{exprel}(-2x)}.
\]
The contribution of \(2h(0)e^{-2x}/(1-e^{-2x})\) beyond the
first interval is exactly
\(-h(0)\log(1-e^{-2h_N})\).
All integration callbacks are meromorphic on their complex
domains; interval division detects their possible poles.

Outward ball arithmetic at \(120\) decimal digits encloses all
integrals, prime terms, and matrix entries. To certify positive
definiteness, an approximate \(LDL^T\) factor supplies an
invertible preconditioner \(R\); the decisive verified inequality is
\[
 \sum_{i,j}|(RGR^T-I)_{ij}|<1.
\]
It implies positive definiteness because the left side bounds
the operator norm of the real symmetric error.
For all six rows it is less than \(2\cdot10^{-94}\).
The trace powers are evaluated after the exact congruence
\((G,B)\mapsto(RGR^T,RBR^T)\), whose first matrix is close
to identity. This controls the error propagation through high
powers. The scalar moments use the ball Taylor series of the
completed zeta function at \(1/2\), followed by its formal
logarithm; no zero ordinate is an input.

The \href{\workbenchurl}{\emph{Arithmetic Weil compression}
program and certificate} give the full enclosures, exact
commands, and verified matrix residuals. The arithmetic uses
Arb's outward interval operations and rigorous integration
\cite{Johansson2017arb,Johansson2018integration}.


\subsection{Removing a certified simple packet}
The coordinate representation also records exactly what remains
after a finite certified simple packet has been removed.

\begin{proposition}\label{prop:finite-simple-deflation}
Let $S$ be a finite set of positive simple zero ordinates, and
write $\pm S$ for its signed counterpart. Delete these coordinates
from the Weil form and from the reciprocal multiplication operator:
\[
 q_S(f,g)=q(f,g)-\sum_{u\in\pm S}
                  \widehat f(u)\overline{\widehat g(u)},
 \qquad C_S=C|_{\{e_u:u\in\pm S\}^{\perp}}.
\]
For $k\ge1$, put $S_k^S=S_k-2\sum_{u\in S}u^{-2k}$. Then
\begin{equation}\label{eq:finite-deflation-excess}
 S_k^S-\operatorname{tr}C_S^{2k}
 =S_k-\operatorname{tr}C^{2k}
 =2\sum_{\gamma>0}(m_\gamma-1)\gamma^{-2k}.
\end{equation}
In the same coefficient basis as $G_N,H_N$, set
$e_u=(\widehat b_1(u),\ldots,\widehat b_d(u))$. The corresponding
matrices satisfy
\begin{align*}
 G_N^S&=G_N-\sum_{u\in\pm S}e_u^*e_u,\qquad
 H_N^S=H_N-\sum_{u\in\pm S}u^{-1}e_u^*e_u,\\
 TH_N^S-G_N^S
   &=(TH_N-G_N)-\sum_{u\in\pm S}(T/u-1)e_u^*e_u.
\end{align*}
At every $T>0$ with $S\subset(0,T)$, both the counted positive zero count and the
distinct positive coordinate count decrease by $|S|$, so their
difference remains $E(T)$.
\end{proposition}

\begin{proof}
In the weighted sampling realization, deleting the coordinates is
the orthogonal projection onto their complement. The image of the
dense test space is dense in that complement: approximate any vector
there in the original space and apply the projection. Thus the
completion of the projected tests, modulo their null space, is
exactly the displayed space for $C_S$. Each deleted coordinate has
weight one because its zero is simple. Its contribution to both
$S_k$ and $\operatorname{tr}C^{2k}$ is $u^{-2k}$. Subtracting
these finitely many terms in the absolutely convergent sums proves
\eqref{eq:finite-deflation-excess}.

For a coefficient column $c$, evaluation is $e_uc$, and its
squared modulus is $c^*e_u^*e_uc$. The reciprocal operator
multiplies that coordinate by $u^{-1}$. These two facts prove
the first two matrix identities; their stated linear combination
proves the third. Deleting a positive simple zero subtracts one
from each of the two positive counts, which proves the last claim.
\end{proof}

Finite deletion can change the approximation error of a chosen
compression. Formula~\eqref{eq:finite-deflation-excess} shows that
the intrinsic multiplicity excess itself is preserved exactly.
Iteration with a height-dependent deleted packet requires an
estimate for the resulting compression error and for the cross
terms with that growing packet. The positive-density argument
above does not supply those estimates for the successive remainders.
